% ---------------------------------------------------------------- титульный лист
\begin{titlepage}
\thispagestyle{empty}
\begin{center}
Министерство образования Республики Беларусь\\[2pt]
Учреждение образования\\[2pt]
<<Белорусский государственный университет\\ информатики и радиоэлектроники>>

\vspace{1.2em}
Факультет компьютерных систем и сетей

\vspace{0.4em}
Кафедра программного обеспечения информационных технологий

\vspace{5em}
А.\,В.~Лебедевич, С.\,Н.~Нестеренков

\vspace{4em}
{\fontsize{17}{22}\selectfont\bfseries ГЕНЕТИЧЕСКИЕ И ЭВОЛЮЦИОННЫЕ ВЫЧИСЛЕНИЯ:\\[2pt]
НЕЙРОННЫЕ СЕТИ И ГЕНЕТИЧЕСКИЕ АЛГОРИТМЫ\par}

\vspace{2em}
Учебно-методическое пособие\\
по дисциплине <<Генетические и эволюционные вычисления>>\\
для студентов II ступени высшего образования

% После решения учебно-методического объединения раскомментировать и заполнить:
% \vspace{1.5em}
% Рекомендовано УМО по образованию в области информатики и радиоэлектроники\\
% в качестве учебно-методического пособия для специальности\\
% <<\ldots>>

\vfill
Минск БГУИР 2026
\end{center}
\end{titlepage}
\setcounter{page}{2}

% ---------------------------------------------------------------- оборот титула
\thispagestyle{empty}
\noindent УДК \rule{4cm}{0.4pt}\\
ББК \rule{4cm}{0.4pt}\\
\hspace*{1.25cm}Л\rule{1cm}{0.4pt}

\vspace{1.5em}
\begin{center}
Р е ц е н з е н т ы:\\[4pt]
\rule{11cm}{0.4pt}\\[6pt]
\rule{11cm}{0.4pt}
\end{center}

\vspace{1.5em}
\noindent\hspace*{1.25cm}\textbf{Лебедевич, А. В.}

\noindent\hangindent=1.25cm\hangafter=0
Генетические и эволюционные вычисления: нейронные сети и генетические алгоритмы : учеб.-метод. пособие /
А.\,В.~Лебедевич, С.\,Н.~Нестеренков. -- Минск : БГУИР, 2026. -- \pageref{LastPage}~с. : ил.

\noindent\hangindent=1.25cm\hangafter=0
ISBN 978-985-\rule{2.2cm}{0.4pt}.

\vspace{1em}
{\small
Пособие содержит теоретический материал и лабораторный практикум по дисциплине <<Генетические и
эволюционные вычисления>>. Рассмотрены основы машинного обучения и искусственных нейронных сетей
(персептрон, линейная сеть, многослойный персептрон и алгоритм обратного распространения ошибки,
самоорганизующиеся карты Кохонена) и генетических алгоритмов (простой генетический алгоритм, теория
схем, операторы селекции, скрещивания, мутации и сокращения популяции, вещественное и перестановочное
кодирование, задачи комбинаторной оптимизации). Практикум включает семь лабораторных работ с
вариантами заданий, указаниями к выполнению, примерами выполнения и контрольными вопросами.

Предназначено для студентов II ступени высшего образования. Может быть использовано для
самостоятельной работы студентов очной и дистанционной форм обучения, а также всеми, кто
изучает методы вычислительного интеллекта и оптимизации.\par}

\vfill
\noindent
\begin{tabular}{@{}p{0.45\textwidth}@{}p{0.55\textwidth}@{}}
\raggedright ISBN 978-985-\rule{2.2cm}{0.4pt} &
\raggedright\arraybackslash\small
\copyright\ Лебедевич А.\,В., Нестеренков С.\,Н., 2026\\
\copyright\ УО <<Белорусский государственный университет информатики и радиоэлектроники>>, 2026
\end{tabular}
\newpage
